\section{Especificaciones de notación, características, escenarios y benchmarks}\label{ap:escenarios}

La Tabla~\ref{tab:notacion} resume los símbolos principales del artículo. La Tabla~\ref{tab:features17} inventaría las $17$ características causales producidas por \texttt{extract\_cf\_features}, indicando su grupo de regularización y si están sujetas a restricción de monotonía no negativa. La Tabla~\ref{tab:config_fixed_hpo} compara las variantes \texttt{SM-CFHR-Fixed} y \texttt{SM-CFHR-HPO}, la Tabla~\ref{tab:escenarios28} enumera los $28$ escenarios oficiales y la Tabla~\ref{tab:benchmarks22} detalla los $22$ procedimientos de referencia.

\begin{table}[!htbp]
\caption{Glosario de notación matemática.}\label{tab:notacion}
\scriptsize
\begin{tabular}{@{}lp{9.5cm}@{}}
\toprule
Símbolo & Significado y unidades \\
\midrule
$M, G$ & Número de plantas ($M=17$) y grupos espaciales ($G=7$) \\
$C_{i,s}, y_{i,s}$ & Capacidad efectiva neta (MW) y generación diaria observada (kWh/día) \\
$E^{\max}_{i,t}$ & Energía nominal máxima diaria al origen: $24\,000\,C_{i,t}$ (kWh/día) \\
$H, H_{\mathrm{dir}}$ & Horizonte en días ($1,7,14,21,28$) y factor de escala directo ($1$ o $H$) \\
$\gamma_h$ & Factor de amortiguamiento geométrico en paso $h$: $0{,}85^{h-1}$ \\
$m_y, \sigma_y$ & Mediana y escala robusta MAD ($1{,}4826\,\MAD$) del objetivo en entrenamiento \\
$\delta, \alpha_{\mathrm{lin}}, \alpha_{\mathrm{meteo}}$ & Umbral de Huber y pesos $L_2$ para dinámica/calendario y meteorología \\
$D_\alpha, \mathcal{M}_+$ & Matriz diagonal de penalización e índices con restricción $w_j \ge 0$ \\
$\mathrm{MAE}_{\mathrm{macro}}$ & Error absoluto medio macro promediado por planta (kWh/día o kWh totales) \\
\bottomrule
\end{tabular}
\end{table}

\begin{table}[!htbp]
\caption{Inventario de características causales de SM-CFHR ($p=7$ sin clima; $p=17$ con clima).}\label{tab:features17}
\scriptsize
\begin{tabular}{@{}cllcc@{}}
\toprule
Índice & Variable en código & Definición & Familia $L_2$ & Monotonía \\
\midrule
1--4 & \texttt{month\_sin/cos}, \texttt{doy\_sin/cos} & Armónicos anuales de calendario al objetivo & $\alpha_{\mathrm{lin}}$ & Libre \\
5 & \texttt{cf30} & Factor de capacidad medio móvil 30d al origen & $\alpha_{\mathrm{lin}}$ & $w_5 \ge 0$ \\
6 & \texttt{diff\_7\_30} & $(\mathrm{cf}_7 - \mathrm{cf}_{30}) \cdot 0{,}85^{h-1}$ & $\alpha_{\mathrm{lin}}$ & Libre \\
7 & \texttt{diff\_1\_7} & $(\mathrm{cf}_1 - \mathrm{cf}_7) \cdot 0{,}85^{h-1}$ & $\alpha_{\mathrm{lin}}$ & Libre \\
8 & \texttt{rad\_pron} & Irradiación horizontal pronosticada (NWP) & $\alpha_{\mathrm{meteo}}$ & $w_8 \ge 0$ \\
9--12 & \texttt{temp/nube/humedad/viento\_pron} & Resto de variables meteorológicas pronosticadas & $\alpha_{\mathrm{meteo}}$ & Libre \\
13--17 & \texttt{anom\_*} & Anomalía causal respecto a media móvil 30d ($\le t_0$) & $\alpha_{\mathrm{meteo}}$ & Libre \\
\bottomrule
\end{tabular}
\end{table}

\begin{table}[!htbp]
\caption{Comparación de configuración entre \texttt{SM-CFHR-Fixed} y \texttt{SM-CFHR-HPO}.}\label{tab:config_fixed_hpo}
\scriptsize
\setlength{\tabcolsep}{3pt}
\begin{tabular}{@{}p{2.3cm}p{4.3cm}p{4.8cm}@{}}
\toprule
Componente & \texttt{SM-CFHR-Fixed} (Primaria) & \texttt{SM-CFHR-HPO} (Secundaria) \\
\midrule
Estructura fija & \multicolumn{2}{p{9.2cm}}{\texttt{n\_knots=0}, \texttt{max\_interactions=0}, \texttt{anchor="none"}, \texttt{structured=True}, \texttt{use\_mono=True}} \\
$\delta$ & $1{,}35$ & TPE en $[0{,}5,\; 3{,}0]$ \\
$\alpha_{\mathrm{lin}}$ & $0{,}01$ & TPE log-uniforme en $[10^{-4},\; 1{,}0]$ \\
$\alpha_{\mathrm{meteo}}$ & $0{,}10$ (clima) / $0{,}01$ (sin clima) & TPE log-uniforme en $[10^{-3},\; 10{,}0]$ (clima) \\
Presupuesto HPO & $B = 0$ evaluaciones & $B = 25$ trials Optuna TPE (3 pliegues, embargo 30d) \\
Optimizador & \multicolumn{2}{p{9.2cm}}{L-BFGS-B (\texttt{max\_iter=400}, \texttt{tol=1e-7}, semilla \texttt{2026})} \\
\bottomrule
\end{tabular}
\end{table}

\begin{table}[!htbp]
\caption{Resumen de las $4$ familias y $28$ escenarios oficiales en \texttt{SCENARIOS\_28}.}\label{tab:escenarios28}
\scriptsize
\begin{tabular}{@{}llccc@{}}
\toprule
Familia & Regímenes & Horizontes $H$ & Clima & $N_{\text{esc}}$ \\
\midrule
$D+1$ (\texttt{d1}) & ET, IP & $1$ día & Con / Sin pronóstico $24\text{ h}$ & $4$ \\
Recursivo sin clima (\texttt{recursive}) & ET, IP & $7, 14, 21, 28$ días & Sin meteorología & $8$ \\
Directo sin clima (\texttt{direct}) & ET, IP & $7, 14, 21, 28$ días & Sin meteorología (acumulado) & $8$ \\
Recursivo + GEFSv12 (\texttt{recursive}) & ET, IP & $7, 14, 21, 28$ días & NOAA GEFSv12 (\texttt{00 UTC}) & $8$ \\
\bottomrule
\end{tabular}
\end{table}

\begin{table}[!htbp]
\caption{Inventario oficial de los $22$ modelos de referencia (excluyendo \texttt{CF-BRH}).}\label{tab:benchmarks22}
\scriptsize
\begin{tabular}{@{}lp{9.2cm}@{}}
\toprule
Categoría & Modelos incluidos \\
\midrule
Bibliotecas ($9$) & \texttt{XGBoost} \cite{chen2016}, \texttt{Random Forest} \cite{breiman2001}, \texttt{Ridge} \cite{hoerl1970}, \texttt{SVR} \cite{drucker1997svr}, \texttt{HistGradientBoosting} \cite{friedman2001gbm}, \texttt{k-NN} \cite{cover1967nearest}, \texttt{Lasso} \cite{tibshirani1996}, \texttt{BayesianRidge} \cite{mackay1992}, \texttt{DecisionTree} \cite{breiman1984cart} \\
Personalizados ($12$) & \texttt{RLAC}, \texttt{IWLAR} \cite{cleveland1979robust,atkeson1997locally}, \texttt{WCFC} \cite{bates1969combination,nikodinoska2022elastic}, \texttt{ABGPR} \cite{rasmussen2006gpml}, \texttt{HARC} \cite{laird1982random,bates2015lme4}, \texttt{SRMOS} \cite{glahn1972mos}, \texttt{ELAC} \cite{brown1959statistical}, \texttt{KLAC} \cite{kalman1960approach,harvey1989forecasting}, \texttt{RARC} \cite{huber1964}, \texttt{RQRF} \cite{koenker1978quantiles,meinshausen2006}, \texttt{PDEC} \cite{hyndman2021fpp3}, \texttt{MARB} \\
Referencia base ($1$) & \texttt{Media móvil de 30 días} (\texttt{MM30}) \\
\bottomrule
\end{tabular}
\end{table}
